\cotefile{147}
\scanpage{133}
$D_{\Delta_r}$ by the $D^*_{\Delta_r}$, to define the
$\Pi_{\Delta_r}$ and $\widetilde\Pi_{\Delta_r}$, \inserted{seems there to}
\deleted{\unsure{will|would}} play a particular role.
\inserted{\unsure{and|here}.}

For \uline{$d_0$ fixed}, the systems
\diagram{\begin{tikzcd}[ampersand replacement=\&]
\widetilde\Pi_{d_0d_1-d_r}
\arrow[r,phantom,"{\unsure{\text{---}|\longrightarrow}}" description]
\&\Pi_{d_0,-d_r}
\end{tikzcd}}
form a homomorphism of semi-
simplicial groupoids, but of an \unsure{essential\textsuperscript{ly}|essentially}
trivial nature (one should howeve\textsuperscript{r} examine
\unsure{on} occasion more closely this so-called
triviality. -- \unsure{?} it can be made explicit without difficulty
in terms of the Galois \unsure{theory|Theory}
of coverings of $D^*_{d_0}$ resp. of $P^*_{d_0}$ \unsure{\ldots|--}).

We shall now introduce
an entirely non-formal element of
structure, which can only be obtained
with hypotheses \unsure{of neighborhoods}
of normal crossings, and expressed
by the fact that the $\widetilde\Pi_{\Delta_r}$ (NB and
not the $\Pi_{\Delta_r}$ !) form a
$1/2$ simplicial groupoid, for $\Delta_r$ variable. In other
\margin{132}
